\section{Introduction}\label{sec:v8-introduction}

Researchers increasingly use large language models to measure long
texts. Once a document exceeds the practical context limit, the
analyst no longer observes the original object: the input to the
final coding step is a compressed surrogate produced by chunking,
retrieval, or recursive summarization. The problem that follows is
one of measurement design as well as prediction. If the compressed
representation drops evidence that depends on combinations across
the document (for example, a tax cut promised in one section paired
with a deficit rule four sections later, where the meaning turns on
the link between the two), the surrogate acts on the wrong object,
and the downstream label inherits that distortion no matter how
accurate the labeler is on its own.

Design-based supervised learning has formalized the analogous
problem one step downstream. \citet{EgamiEtAl2024} show that even
highly accurate surrogate labels can bias downstream estimates when
they substitute for gold-standard labels without appropriate
sampling and overlap conditions. Their fix is to log a sampling
design over the trusted labels and use it to make the downstream
inference valid. Long-document measurement adds an earlier failure
mode: the surrogate is both a compressed representation of the
original document and a predicted label. A retrieval or
long-context score may be accurate on average and still tell the
analyst nothing about which cross-document interactions survived
compression.

C-TreePO addresses the representation-level problem. A compression
tree partitions a document into contiguous leaves, summarizes each
leaf into an evidence state, merges sibling states up the tree, and
scores the root. The state map \(g\) is valid for an oracle
\(f^\ast\) when three local laws hold on the realized tree.
Sufficiency: leaf summaries preserve the evidence relevant to the
task. Idempotence: a state stays invariant when passed through the
summarizer again. Merge consistency: a merged state preserves the
parent span it represents. The laws are local in scope and global
in consequence: a valid root state can replace the original document
for any oracle-compatible readout.

The C-TreePO objective has the same two-level shape. The root
channel fits the observed document-level expert target. The
local-law channel trains and audits the state map inside the
realized tree. Direct \(f^\ast\) calls are expensive, so we use a
fitted proxy \(\widehat f\) for the dense local supervision and
sample \(f^\ast\) at a small set of nodes to correct the
proxy/oracle gap.

Audit happens at tree nodes, so the analyst picks the unit: a leaf,
a paragraph, a section, or the full document. Design-based
machinery aggregates the sampled node labels into a document-level
guarantee, and the same design trades one expensive full-document
expert read for several cheaper reads of shorter units at
comparable accuracy. The same supervision design also trains the
compression operators under a fixed audit budget; that is the
optimization side of the framework, the ``PO'' in C-TreePO.

The certificate makes the resulting measurement error reportable.
Section~\ref{sec:v8-audit-certificate} decomposes the bound on
downstream estimate drift into four interpretable terms. A sampled
local-distortion term records how far the produced states moved
from the raw spans they represent on the realized tree. A
calibration term tracks the gap between the accessible local judge
and the target oracle: the part of error that more sampled labels
cannot remove. A sampling term carries Horvitz--Thompson uncertainty
from the audit design. A clipping term prices any variance control
applied to inverse-propensity weights. Each term points at a
specific lever: train the state map harder, change the local judge,
sample more nodes, or report the clipping bias explicitly.

We instantiate the framework on the Manifesto LLM benchmark of
\citet{BenoitEtAl2025}. The Comparative Manifestos Project (CMP)
has collected and hand-coded party election platforms for decades;
its expert-survey arm produces mean placements of each platform on
seven-point policy scales. \citet{BenoitEtAl2025} introduce an LLM
scoring pipeline for that target: it summarizes each
manifesto--dimension pair, then scores the summary against the
expert mean on six policy dimensions, from economic policy to
decentralization. We rerun the same task with C-TreePO on the same
\(235\) party platforms (hereafter, the \emph{Manifesto LLM
replication}).
Per-dimension trees reach macro \(r=0.829\); the universal-summary
tree reaches \(0.807\). Both clear the matched open-weight baseline
at \(0.793\), and the per-dimension tree clears the proprietary
18-score ensemble at \(0.817\).

Two controlled checks build intuition for the mechanism behind the
application. The Markov check uses a known oracle and a known
sufficient state to show when local labels can substitute for some
root labels under a fixed token budget. The IPW check uses hard
synthetic audit designs to show how logged-propensity correction
keeps the node-level estimand centered while variance grows in the
expected places.

The framework, the application, and the controlled checks share one
object: a learned state map whose nodes satisfy three local laws on
the realized tree. We define that object and the laws next, write
the training objective and the audit, run the controlled checks,
and turn to the manifesto application.
